\section[Rédaction Bourbaki n\textsuperscript{o}\,457 : degré et nombre de points]{Rédaction Bourbaki n\textsuperscript{o}\,457 : degré et nombre de points d'une algèbre séparable\protect\verifie{Bourbaki457}}\label{sec:bourbaki-457}

Cette section ne porte pas sur une cote du fonds de Montpellier. La rédaction n\textsuperscript{o}\,457 des archives Bourbaki, datée d'octobre 1965 par sa notice (Archives Henri Poincaré, \href{https://archives-bourbaki.ahp-numerique.fr/items/show/792}{item 792}), qui la donne à Grothendieck, est un projet de réédition du chapitre~V d'\emph{Algèbre}, consacré aux corps commutatifs. Elle y substitue l'étude des algèbres à celle des extensions, et classe les algèbres étales sur un corps par les ensembles finis sur lesquels opère le groupe de Galois. Elle est placée dans cette partie pour cette théorie de Galois des algèbres étales. Sa transcription reste locale pour l'instant ; on n'en cite ici aucun passage.

\emph{Les points d'une algèbre.} Soient $k$ un corps et $\Omega$ une extension algébriquement close de $k$. Pour une $k$-algèbre commutative $A$, on note $P(A)$ l'ensemble des homomorphismes de $k$-algèbres de $A$ dans $\Omega$. En langage géométrique, ce sont les points de $\operatorname{Spec} A$ à valeurs dans $\Omega$. Pour $A = k[X]/(f)$, un homomorphisme $A \to \Omega$ est déterminé par l'image de $X$, qui doit être une racine de $f$ : les points de $A$ sont les racines de $f$ dans $\Omega$. Pour $k = \mathbf Q$ et $f = X^2 - 2$, on trouve les deux plongements de $\mathbf Q(\sqrt 2)$ dans $\mathbf C$.

\emph{En degré fini, la dimension compte les points.} Si $f$ est séparable de degré $n$, il a $n$ racines distinctes dans $\Omega$, et $A = k[X]/(f)$ a exactement $n$ points pour une dimension $n$. Plus généralement, une algèbre étale de degré $n$ devient, après extension des scalaires à $\Omega$, le produit $\Omega^n$, dont les $n$ projections sont les points. La séparabilité est indispensable : pour $k$ de caractéristique $p$ et $a \in k$ qui n'est pas une puissance $p$-ième, l'algèbre $k[X]/(X^p - a)$ est de dimension $p$ et n'a qu'un point, car $X^p - a$ n'a qu'une racine dans $\Omega$.

\emph{La page.} Le corollaire~6 de la proposition~14 (PDF, p.~32--33) considère une algèbre $A$ entière et séparable sur $k$, c'est-à-dire réunion filtrante de sous-algèbres étales de degré fini $A_i$, et affirme que son degré $[A:k]$ est égal au cardinal de $P(A)$ ; en particulier, $A$ est de degré fini si et seulement si $P(A)$ est fini. La démonstration ne traite que ce dernier point. Les $P(A_i)$ sont finis, $P(A)$ est leur limite projective, et les applications de transition sont surjectives ; un lemme sur les limites projectives d'ensembles finis non vides donne alors des surjections de $P(A)$ sur chaque $P(A_i)$, dont les cardinaux tendent vers l'infini. On a donc montré que $P(A)$ est infini lorsque $A$ est de degré infini, et rien de plus. Une note dactylographiée le reconnaît : le rédacteur y constate qu'il a donné à \q{card} la seule valeur $\infty$ pour tout ensemble infini. Un lecteur a écrit à l'encre, à droite de la formule, que l'égalité est fausse en dimension infinie, avec les mots \q{dénombrable} et \q{continu}. Le contre-exemple qui suit rend précise cette objection.

\emph{Pourquoi l'égalité échoue en degré infini.} Une base dénombrable ne voit que des combinaisons finies, tandis qu'un point de $A$ est un choix compatible à tous les étages $A_i$ à la fois. L'ensemble de Cantor $C = \{0, 1\}^{\mathbf N}$ montre l'écart. Ses parties à la fois ouvertes et fermées sont les réunions finies de cylindres, un cylindre étant l'ensemble des suites qui commencent par un mot fini donné ; il n'y en a qu'une infinité dénombrable. Ses points sont les suites infinies de $0$ et de $1$, qui ont la puissance du continu. Les fonctions localement constantes $C \to k$ sont les combinaisons finies d'indicatrices de cylindres : elles forment une algèbre de dimension dénombrable, et chaque point de $C$ en donne un point, l'évaluation en ce point. Étage par étage, les fonctions qui ne dépendent que des $n$ premiers termes forment l'algèbre étale $k^{\{0,1\}^n}$, de degré $2^n$ et à $2^n$ points ; chaque étage double les deux nombres, mais à la limite la dimension reste dénombrable et les points deviennent les suites infinies (figure~\ref{fig:457-cantor}).

\begin{figure}[htbp]
\centering
\begin{tikzpicture}[x=0.62cm, y=0.95cm, >=stealth, every node/.style={font=\small}]
% le cylindre [01] : les branches qui passent par le noeud 01
\fill[piste!12] (-2,-2) -- (-3.95,-4.75) -- (-0.05,-4.75) -- cycle;
\foreach \d in {0,...,3} {
  \pgfmathtruncatemacro{\n}{2^\d - 1}
  \foreach \i in {0,...,\n} {
    \pgfmathsetmacro{\x}{(\i + 0.5)*2^(4 - \d) - 8}
    \pgfmathsetmacro{\xl}{(2*\i + 0.5)*2^(3 - \d) - 8}
    \pgfmathsetmacro{\xr}{(2*\i + 1.5)*2^(3 - \d) - 8}
    \draw[black!55] (\x,-\d) -- (\xl,-\d-1);
    \draw[black!55] (\x,-\d) -- (\xr,-\d-1);
    \fill (\x,-\d) circle (1.6pt);
  }
}
\foreach \i in {0,...,15} { \fill (\i + 0.5 - 8,-4) circle (1.6pt); \draw[black!40, dotted] (\i + 0.5 - 8,-4) -- (\i + 0.5 - 8,-4.6); }
\node[above left=-1pt] at (-2,-0.5) {\footnotesize $0$};
\node[above right=-1pt] at (2,-0.5) {\footnotesize $1$};
% la branche 1011...
\draw[frise, very thick] (0,0) -- (4,-1) -- (2,-2) -- (3,-3) -- (3.5,-4);
\draw[frise, very thick, dotted] (3.5,-4) -- (3.5,-4.6);
\node[frise, below, align=center] at (3.5,-4.6) {\footnotesize branche $1011\ldots$,\\\footnotesize un point de $C$};
\node[piste, below, align=center] at (-2,-4.6) {\footnotesize cylindre des suites\\\footnotesize qui commencent par $01$};
\fill[piste] (-2,-2) circle (2.2pt);
\node[piste, above left] at (-2,-2) {\footnotesize $01$};
\foreach \d in {0,...,4} { \pgfmathtruncatemacro{\p}{2^\d} \node[black!65, anchor=west] at (8.4,-\d) {\footnotesize $2^{\d} = \p$}; }
\end{tikzpicture}
\caption{L'arbre binaire des mots finis, jusqu'à la longueur~$4$. Chaque n\oe ud est un mot $v$ et donne le cylindre des suites qui commencent par $v$ : il y en a $2^n$ à l'étage~$n$, une infinité dénombrable en tout, et leurs indicatrices engendrent l'algèbre $A$. Un point de $C$ est une branche infinie de l'arbre, un choix à chaque étage : il y en a la puissance du continu. La branche dessinée appartient aux cylindres $1$, $10$, $101$, $1011$, et non au cylindre $01$.}
\label{fig:457-cantor}
\end{figure}

\begin{contrexemple}[Bourbaki, rédaction n\textsuperscript{o}\,457, corollaire~6]
Soient $k$ un corps, $C = \{0, 1\}^{\mathbf N}$ et $A$ la sous-$k$-algèbre de $k^C$ formée des fonctions qui ne dépendent que d'un nombre fini de termes de la suite. Tout élément de $A$ est racine d'un polynôme unitaire séparable scindé sur $k$ ; en particulier $A$ est entière et séparable sur $k$. On a $[A:k] \leq \aleph_0$, et, pour toute $k$-algèbre non nulle $\Omega$, en particulier toute extension algébriquement close de $k$, l'ensemble des homomorphismes de $k$-algèbres $A \to \Omega$ a au moins la puissance du continu. Donc $[A:k] < \operatorname{card} P(A)$.
\end{contrexemple}

\begin{proof}
Si $f$ ne dépend que des $n$ premiers termes, elle ne prend qu'un nombre fini de valeurs $c_1, \dots, c_r$, deux à deux distinctes, et $f$ est racine de $\prod_j (X - c_j)$, qui est unitaire, scindé sur $k$ et à racines simples, donc séparable : au point $z$, le facteur $f(z) - c_j$ correspondant à la valeur $f(z)$ est nul. Pour chaque mot $v$ de longueur $n$, soit $\chi_v$ l'indicatrice du cylindre des suites qui commencent par $v$. Alors $f = \sum_v f(\bar v)\, \chi_v$, où $\bar v$ est la suite $v$ prolongée par des zéros : en un point $z$, seul le terme du mot formé des $n$ premiers termes de $z$ subsiste, et $f(\bar v) = f(z)$ puisque $\bar v$ et $z$ ont les mêmes $n$ premiers termes. Les $\chi_v$, pour tous les mots finis $v$, forment une famille dénombrable qui engendre $A$, d'où $[A:k] \leq \aleph_0$. Pour un point $x$ de $C$, l'évaluation $f \mapsto f(x)$, suivie de l'application $k \to \Omega$, est un homomorphisme de $k$-algèbres $A \to \Omega$. Si $x \neq y$, ils diffèrent en un rang $m$, et la fonction qui vaut $1$ sur les suites de terme $1$ au rang $m$ et $0$ ailleurs appartient à $A$ et prend les valeurs $1$ et $0$ en $x$ et $y$, distinctes dans $\Omega$ non nulle. L'application $x \mapsto \mathrm{ev}_x$ est donc injective, et $\operatorname{card} P(A) \geq \operatorname{card} C = 2^{\aleph_0} > \aleph_0$.
\end{proof}

\emph{Vérification en Lean.} Lean vérifie l'énoncé ci-dessus : $A$ y est définie directement comme l'algèbre des fonctions qui ne dépendent que d'un nombre fini de termes, et l'énoncé final dit que, pour toute extension algébriquement close $\Omega$ de $k$, il existe une $k$-algèbre entière dont tout élément est racine d'un polynôme séparable scindé, et dont le degré diffère du cardinal de l'ensemble de ses homomorphismes dans $\Omega$. Trois points ne sont pas formalisés : l'identification de $A$ avec les fonctions localement constantes sur l'ensemble de Cantor, les égalités $[A:k] = \aleph_0$ et $P(A) \simeq C$ (pour $\Omega$ un corps), dont la réfutation n'a pas besoin, et le cas du degré fini, où le corollaire est juste, ainsi que sa dernière assertion, qui est juste aussi.

\begin{remarque}
Le même phénomène se rencontre pour les corps. L'extension $K = \mathbf Q(\sqrt 2, \sqrt 3, \sqrt 5, \dots)$ engendrée par les racines carrées des nombres premiers est de degré dénombrable sur $\mathbf Q$, et ses plongements dans $\mathbf C$ correspondent aux éléments de son groupe de Galois $(\mathbf Z/2\mathbf Z)^{\mathbf N}$, qui a la puissance du continu. En degré infini, le cardinal de $P(A)$ n'est d'ailleurs pas fixé par le degré : pour un ensemble infini $D$, l'algèbre des fonctions $D \to k$ constantes hors d'une partie finie a pour degré le cardinal de $D$, et ses points, les évaluations en les éléments de $D$ et la valeur constante à l'infini, sont en même nombre. Ces deux exemples ne sont pas formalisés. La forme juste du corollaire est sa dernière assertion, que la démonstration de la page établit.
\end{remarque}
